\centerline{\Huge \bf  Тренировочная Олимпиада 3}
\centerline{\Huge \bf 8 класс}
\centerline{\Huge \bf Уровень: Регион ВсОШ. День 2}


\begin{flushright}
    \scriptsize{Глава 3. День, когда Токио-3 замер.}
\end{flushright}


\problem{1}
На день рождения Расула Владимировича пришло 49 учеников ЭлМШ. Известно, что любые 32 ребёнка суммарно съели хотя бы половину от всей еды. Какую наибольшую часть от еды смог бы съесть один Эрдем? Расул Владимирович к еде не прикасался, а дети по итогу съели всю имеющуюся еду.
\begin{flushright}
    \textit{\small{(Гасанов Р.В.)}}
\end{flushright}

\problem{2}     
Сколькими способами, меняя между собой некоторые буквы (не обязательно различные), из слова БУСЯДУСЯЛЮСЯМАРУСЯКРИСТИВАНГА можно получить другое слово, которое тоже означало бы все имена кошек Расула Владимировича.
\begin{flushright}
    \textit{\small{(Гасанов Р.В.)}}
\end{flushright}

\problem{3}
В треугольнике $ABC$ угол $B$ равен $60^\circ$. Точка $O \ -$ центр описанной окружности $\omega$, точка $H \ -$ ортоцентр, а точка $M \ -$ середина меньшей дуги $AC.$ Докажите, что $BM \perp OH.$  
\begin{flushright}
    \textit{\small{(Гасанов Р.В.)}}
\end{flushright}

\problem{4} 
На плоскости нарисован график параболы $y = x^2$. К ней проведены две касательные $l_1, l_2$, касающиеся в точках $M$ и $N$, пересекающиеся в точке $O$ под углом $90^{\circ}$ и точка $O$ находится на продолжении оси симметрии параболы. Окружность $\omega$ касается параболы в её вершине, а так же  
прямых $l_1$ и $l_2$ в точках $K, L$ соответственно. Найдите площадь четырёхугольника $KLNM.$
\begin{flushright}
    \textit{\small{(Гасанов Р.В.)}}
\end{flushright}

\problem{5}
Назовём \textit{паспортом} натурального числа набор простых чисел, входящих в его разложение. Докажите, что в любой натуральной арифметической прогресии есть сколь угодно много чисел с одинаковым паспортом.
\begin{flushright}
    \textit{\small{(Гасанов Р.В.)}}
\end{flushright}

%У красноухого бабуина Вячеслава есть $N$ книг Кормана по алгоритмам. Однажды Вячеслав забрался на 163-этажное здание Бурдж-Халифа и начал скидывать оттуда книги Кормана с разных этаже. Если книга не разрывалась при приземлении, то он подбирал её и мог скинуть ещё раз с какого-нибудь этажа Бурдж-Халифы, а если всё же разорвалась, то увы скидывать он её больше не сможет. За какое минимальное количество бросков бабуин сможет наверняка понять, начиная с какого этажа книга по алгоритмам полностью разрывается?